\chapter{1994:George A. Olah}

\label{chap:chemistry1994}

The Nobel Prize in Chemistry for the year 1994 was awarded to George A. Olah.

\textbf{Motivation:}

for his contribution to carbocation chemistry

\section{George A. Olah}

\subsection*{Early Life and Education}

George A. Olah was born in 1927 in Budapest, Hungary.

\subsection*{Career and Major Research}

Olah earned his Ph.D. from the Technical University of Budapest in 1949. After leaving Hungary in 1956, he worked at Dow Chemical in Canada and then at Case Western Reserve University. From 1977 he was a professor at the University of Southern California, where he founded the Loker Hydrocarbon Research Institute.

\subsection*{Nobel Prize-winning Work}

The work for which George A. Olah was awarded the Nobel Prize in Chemistry in 1994 was described by the Nobel Committee as: "for his contribution to carbocation chemistry".

He used superacids to prepare stable, long-lived carbocations, and he showed that carbocations are real intermediates in many chemical reactions, resolving a long-standing debate in organic chemistry.

\subsection*{Significance and Impact}

His work underpinned understanding of reactions that are central to petroleum chemistry and refined synthetic organic chemistry.

\subsection*{Later Career and Honors}

Olah remained at the University of Southern California. He received the 1994 Nobel Prize in Chemistry, and he died in 2017.

\subsection*{Legacy}

The study of carbocations is now fundamental to organic chemistry, and Olah's later concept of a methanol economy influenced energy research.

\subsection*{Selected Publications}

"My Search for Carbocations and Their Role in Chemistry" (Angewandte Chemie International Edition in English, 1995)

\clearpage

\section*{Chapter Summary}

The 1994 Nobel Prize in Chemistry recognized contributions to carbocation chemistry. Olah demonstrated that carbocations can be prepared and studied as stable species in superacid media, establishing their role as reaction intermediates.
